% Generated draft: requires mathematical and source review.
\documentclass[11pt]{article}
\usepackage[margin=25mm]{geometry}
\usepackage{fontspec}
\setmainfont{Latin Modern Roman}
\usepackage{amsmath,amssymb,mathtools,unicode-math}
\setmathfont{Latin Modern Math}
\usepackage{enumitem,needspace,xurl,hyperref}
\setlength{\parindent}{0pt}
\setlength{\parskip}{5pt}
\setlength{\emergencystretch}{3em}
\newcommand{\N}{\mathbb{N}}
\newcommand{\Z}{\mathbb{Z}}
\newcommand{\Q}{\mathbb{Q}}
\newcommand{\R}{\mathbb{R}}
\newcommand{\C}{\mathbb{C}}
\newcommand{\F}{\mathbb{F}}
\newcommand{\A}{\mathbb{A}}
\newcommand{\PP}{\mathbb{P}}
\newcommand{\E}{\mathbb{E}}
\newcommand{\eps}{\varepsilon}
\newcommand{\dd}{\,\mathrm{d}}
\begin{document}
\section*{UnsolvedMath ID: 30003589}
\textbf{Source problem code:} catalogue entry
\subsection{Definitions and mathematical statement}
The supplied record combines two separate questions and is marked \texttt{statement\_status: unrecoverable}. The formulations below retain both components; they do not constitute a recovered single problem.

\begin{enumerate}
\item An approximately finite-dimensional (AF) \(C^*\)-algebra is understood through approximation of finite subsets in norm by finite-dimensional \(C^*\)-subalgebras. The source does not explicitly impose separability; the usual sequential inductive-limit convention would impose it. A uniformly hyperfinite (UHF) algebra, under the usual convention, is a unital inductive limit of full matrix algebras with unital connecting maps.

For a positive integer \(n\), put \(\mathbb Z_n=\mathbb Z/n\mathbb Z\). An action is a homomorphism \(\alpha\colon\mathbb Z_n\to\operatorname{Aut}(A)\), and \(A\rtimes_\alpha\mathbb Z_n\) denotes its crossed product. No simplicity, outerness, or freeness hypothesis is supplied. For this finite group, the full and reduced crossed products agree.

In the usual separable setting, a \(C^*\)-algebra \(B\) satisfies the universal coefficient theorem (UCT) if, for every separable \(C^*\)-algebra \(D\), the natural map from Kasparov theory to homomorphisms of \(K\)-groups occurs in the UCT exact sequence
\[
\begin{split}
0\longrightarrow{}&\operatorname{Ext}^1_{\mathbb Z}(K_0(B),K_1(D))
\oplus\operatorname{Ext}^1_{\mathbb Z}(K_1(B),K_0(D))\\
\longrightarrow{}&KK_0(B,D)\\
\longrightarrow{}&\operatorname{Hom}_{\mathbb Z}(K_0(B),K_0(D))
\oplus\operatorname{Hom}_{\mathbb Z}(K_1(B),K_1(D))
\longrightarrow0.
\end{split}
\]
Here \(K_i\) denotes operator \(K\)-theory and \(KK_0\) denotes the even Kasparov group.

With the separable convention explicitly imposed, the inheritance question is whether
\[
\forall A\text{ separable AF},\quad
\forall n\geq1,\quad
\forall\alpha\colon\mathbb Z_n\to\operatorname{Aut}(A),\qquad
A\rtimes_\alpha\mathbb Z_n\text{ satisfies the UCT}.
\]
The source also asks whether the general UCT problem can be reduced to these crossed products, and whether the answers change upon restricting \(A\) to UHF algebras. Under the customary meaning of the general UCT problem, one possible precise reduction target is the implication
\[
\begin{gathered}
\text{Every finite cyclic crossed product of a separable AF algebra satisfies the UCT}\\
\Longrightarrow\\
\text{Every separable nuclear \(C^*\)-algebra satisfies the UCT}.
\end{gathered}
\]
This implication is an interpretation of the requested reduction, not an assertion that it holds. The supplied record does not specify the intended reduction or whether its UHF question concerns inheritance, reduction, or both. Without resolving the separability convention, the displayed inheritance statement is only a conditional formulation of the source question.

\item Separately, the source asks: is the hyperfinite factor \(R\) of type \(\mathrm{II}_1\) quasidiagonal? Under the standard interpretation, \(R\) is the hyperfinite finite von Neumann factor with separable predual, equipped with its normalized trace.

To distinguish possible meanings, a representation \(\pi\colon B\to\mathcal B(H)\) is quasidiagonal if there are finite-rank orthogonal projections \(P_\lambda\) such that
\[
P_\lambda\longrightarrow I_H\text{ strongly},\qquad
\|P_\lambda\pi(b)-\pi(b)P_\lambda\|\longrightarrow0
\quad(b\in B).
\]
For separable \(H\), this is conventionally expressed using a sequence. Quasidiagonality of an abstract \(C^*\)-algebra concerns the existence of a faithful quasidiagonal representation, rather than a prescribed representation. The record does not specify whether it intends this property for the underlying, norm-nonseparable \(C^*\)-algebra of \(R\), quasidiagonality in its standard representation, or a particular extension of the separable definition. These interpretations must be resolved before promoting a precise standalone statement.
\end{enumerate}
\subsection{Short English statement}
The record combines a UCT question for finite cyclic crossed products of AF algebras, including reduction and UHF variants, with a separate question about quasidiagonality of the hyperfinite \(\mathrm{II}_1\) factor. Separability, the intended reduction, and the meaning of quasidiagonality for \(R\) require clarification. The supplied open label is dated curation, not an independently verified conclusion.
\subsection{Sources}
\begin{itemize}
\item Local UnsolvedMath registry: \texttt{data/unsolvedmath/problems.json}, permanent ID \(30003589\), source code \texttt{catalogue entry}; supplied title: ``Composite of Unrelated Operator-Algebra Questions.'' No individual proposer is supplied; the proposed year is \(2017\).
\item Supplied workshop citation: \emph{Mini-Workshop: MASAs and Automorphisms of C*-Algebras (2018)}, Oberwolfach Reports, DOI \texttt{10.4171/owr/2017/42}, pp.\ 2601--2629. \url{https://doi.org/10.4171/owr/2017/42}. This is the supplied primary-source attribution; its contents have not been independently checked here.
\item Supplied literature reference: Barlak--Szabo, ``crossed products of the Razak--Jacelon algebra.'' \url{https://arxiv.org/abs/1712.00823}.
\item Supplied literature reference: Barlak--Li, ``Cartan subalgebras and the UCT problem.'' \url{https://arxiv.org/abs/1511.02697}.
\item The supplied literature assessment was checked at \texttt{2026-08-21}. This is supplied dated curation, not independently verified. It labels the problem open, describes the crossed-product inheritance question as tied to the unresolved UCT problem, attributes to Barlak--Szabo a reduction to special crossed products of the Razak--Jacelon algebra, and attributes to Barlak--Li permanence under additional Cartan hypotheses. The background explicitly requests independent verification of time-sensitive status and cited claims.
\item The supplied statement-verification note reports that the official problem-session list separates the AF crossed-product UCT question from the hyperfinite-factor quasidiagonality question. The record retains the composite extraction for provenance, supplies no clean statement, and marks its statement status unrecoverable. This verification note is itself supplied evidence, not a claim of source consultation here.
\end{itemize}
\end{document}
